\BourbakiStarSection{历史注记}

（方括号中的数字指本注末尾的参考文献。）

数学中几乎没有比合成律更原始的概念了：它似乎与自然数和可测量量的最初计算基础密不可分。流传至今的关于埃及人和巴比伦人数学的最古老文献表明，他们早已拥有了一套关于自然数 \textgreater0 、有理数 \textgreater0 、长度与面积的完整的计算规则体系；尽管保存下来的文本仅涉及给定数量具有明确数值的问题，† 但它们使我们毫不怀疑他们对所用规则所赋予的普遍性，并展示了在处理一次和二次方程方面相当卓越的技术能力（{[}1{]}，第 179 页及以下）。另一方面，完全没有丝毫关注去证明所使用的规则，甚至对所出现的运算也没有精确的定义：后者与前者一样，都是以纯粹经验的方式产生的。

然而，正是这样一种关注，在古典时代的希腊人的著作中已经非常清晰地显现出来；诚然，自然数理论并未找到公理化的处理（这样的公理化直到十九世纪末才出现；参见《集合论》第 IV 章的历史注记）；但在欧几里得的《几何原本》中，有许多段落对计算规则给出了形式化的证明，这些规则与整数计算的规则一样，在直觉上是“显然”的（例如，两个有理数乘积的交换性）。这类性质最引人注目的证明，是关于量理论的证明，这是希腊数学最具原创性的创造（众所周知，它相当于我们的实数理论 \textgreater0 ；参见《一般拓扑学》第 IV 章的历史注记）；在这里，欧几里得除其他事项外，考虑了两种量之比（ratio）的乘积，并表明该乘积与这些比的表示形式无关（这是 § 1 第 6 号中按等价关系对合成律取“商”的第一个例子），并且这些乘积满足交换律（{[}2{]}，第五卷，命题 22—23）。†

然而，不可隐瞒的是，这种向严格性的进步在欧几里得那里伴随着停滞，甚至在代数计算技术方面在某些方面还有所倒退。几何的压倒性优势（显然是在其光照下构思了量论的）使得代数符号的任何独立发展都陷入瘫痪：进入计算的元素必须始终在几何上被“表示”；此外，出现的两种合成律并非定义在同一集合上（比值的加法并未以一般方式定义，而两条长度的乘积不是长度而是面积）；其结果是缺乏灵活性，使得操纵次数大于二的代数关系几乎不可行。

只有当古典希腊数学衰落时，丢番图才回归到“计算师”或专业计算者的传统，这些人一直沿用从埃及人和巴比伦人那里继承来的规则：他不再用几何表示来束缚自己所考虑的“数”，自然便发展出抽象代数计算的规则；例如，他给出的规则（用现代语言来说）对 m 和 n 的小值（正或负）等价于公式 \(x^{m+n} = x^{m}x^{n}\) （{[}3{]}，第 I 卷，第 8—13 页）；稍后（第 12—13 页），陈述了“符号法则”，即负数计算的初步内容‡；最后，丢番图首次使用一个字母来表示方程中的未知数。相比之下，他似乎很少关心为他解决问题所应用的方法赋予普遍意义；至于由欧几里得开创的合成律的公理化观念，这对丢番图及其直接后继者的思想而言似乎是陌生的；这种观念直到 19 世纪初才在代数学中重新出现。

在随后的几个世纪中，首先需要发展一套足以表达抽象规律的代数符号体系，另一方面则需要通过观察足够多样的特例，将“数”的概念加以拓展，从而形成一般性的观念。为此，希腊人创立的量之比的公理化理论并不足够，因为它只不过使实数 \textgreater0 的直观概念以及这些数上的运算变得精确而已，而这些运算巴比伦人早已以较为混乱的形式知晓；如今所考虑的“数”将是希腊人从未有过的概念，并且起初也没有可感知的“表示”：一方面是零和负数，它们出现在中世纪盛期的印度数学中；另一方面是虚数，这是 16 世纪意大利代数学家的创造。

除了零——最初是作为分隔符号引入，后来才被视为数（见《集合论》第 III 章的历史注记）——这些扩充的共同特点是（至少起初）它们纯粹是“形式的”。这里必须理解为，新的“数”首先表现为对某些情形进行运算的结果，而在这些情形中，若严格遵循其定义，这些运算本无意义（例如，两个自然数之差 \(a - b\) 当 \(a < b\) 时的情形）：由此产生了赋予它们的名称：“虚假的”、“虚构的”、“荒谬的”、“不可能的”、“想象的”数等等。对于古典时期的希腊人而言，他们尤其热爱清晰的思想，这样的扩展是不可想象的；它们只能出现在那些比希腊人更倾向于对其方法的力量表现出某种神秘信仰（如 18 世纪所说的“分析的普遍性”），并任由自己随计算的机械过程前行而不去探究每一步是否有充分依据的计算者手中；这种信心事后通常被证明是合理的，因为把这些严格说来仅对已知数有效的计算规则应用于这些新的数学实体，导出了精确的结果。这就解释了为何这些数字概念的推广——起初仅作为一系列运算中的中介出现，其起点和终点都是真正的“数字”——逐渐被独立看待（不依赖于具体计算的应用）；一旦迈出这一步，人们便寻求或多或少可感知的解释，从而赋予新对象在数学中存在的权利。†

在这个问题上，印度人已经意识到在某些情况下负数应被赋予的解释（例如商业问题中的债务）。在随后的几个世纪里，随着希腊和印度数学的方法与成果（经由阿拉伯人）传播到西方，这些数的运算变得更加熟悉，并获得了其他具有几何或运动学性质的“表示”。而代数符号体系的逐步改进，则是中世纪末期代数学中唯一显著的进步。

16 世纪初，由于意大利学派数学家发现了三次方程及随后四次方程的“根式”解法（我们将在《代数》第 V 章的历史注记中更详细地讨论这一点），代数学获得了新的推动力；正是在此时，尽管他们心存厌恶，却不得不将虚数引入计算之中；此外，对这些“不可能”数的运算信心逐渐建立起来，正如对负数的运算一样，尽管在两个多世纪里，人们从未构想出它们的任何“表示”。

另一方面，代数符号体系由韦达和笛卡尔决定性地完善；自笛卡尔之后，代数符号大体上就是我们今天所使用的形式。

从 17 世纪中叶到 18 世纪末，微积分的创立所开辟的广阔视野似乎导致了代数学总体上的一定忽视，尤其是对合成律以及实数与复数性质的数学反思。† 因此，自 17 世纪末在力学中已广为人知的力的合成与速度的合成，在代数学中并未产生反响，尽管它们已经蕴含了向量演算的萌芽。事实上，必须等到大约 1800 年前后，导致复数几何表示的思想运动（见《一般拓扑学》第 VIII 章的历史注记），才看到向量加法在纯数学中得到应用。‡

大约在同一时期，代数学中首次将合成律的概念沿两个不同方向推广到与“数”（迄今为止该词最广义的含义）仅有遥远类比关系的元素上。这两个推广中的第一个归功于 C. F. 高斯，起因于他对二次型 \(ax^{2} + bxy + cy^{2}\) （具有整数系数）的算术研究。拉格朗日曾在具有相同判别式的形式集合上定义了一个等价关系†，并且另一方面证明了一个恒等式，该恒等式为该集合提供了一个交换合成律（并非处处有定义）；基于这些结果，高斯证明了该合成律与上述等价关系的一种细化形式相容（在 § 1 第 6 号的意义下）（{[}5{]}，第 I 卷，第 272 页）：“这表明”，他随后说道，“必须理解两个或若干类的复合类是什么意思”。接着他继续研究“商”律，实质上确立了它（用现代语言说）是一个交换群律，并且所用的论证其普遍性在很大程度上远远超出了高斯所考虑的特殊情形（例如，他证明逆元唯一性的论证与我们针对 § 2 第 3 号中任意合成律所给出的论证完全相同（同上，第 273 页））。他并未止步于此：稍后回到这个问题时，他认识到类的复合与模素数整数的乘法之间的类比‡（同上，第 371 页），但也证明了具有给定判别式的二次形式的类群并不总是循环群；他在这方面给出的指示清楚地表明，他至少在这个特殊情形下已经认识到了有限阿贝尔群的一般结构，我们将在第七章中研究这一结构（{[}5{]}，第 I 卷，第 374 页及第 II 卷，第 266 页）。

我们想谈的另一系列研究也以群的概念及其最终引入数学而告终：这就是“代换理论”，即拉格朗日、范德蒙德和高斯关于代数方程求解思想的发展。我们在此不给出这一主题的详细历史（参见《代数》第 V 章的历史注记）；我们必须提及鲁菲尼以及随后柯西（{[}6{]}，（2），第 I 卷，第 64 页）对有限集合的两个置换的“乘积”的定义§，以及关于有限置换群的初步概念：传递性、本原性、单位元、可交换元素等。但这些初步结果总体上仍然相当肤浅，埃瓦里斯特·伽罗瓦必须被视为该理论的真正开创者：在他那篇令人难忘的著作 {[}8{]} 中，他将代数方程的研究归结为与之相关的置换群的研究，从而大大深化了对后者的研究，无论是在群的一般性质方面（伽罗瓦是第一个定义正规子群的概念并认识到其重要性的人），还是在具有特殊性质的群的确定方面（他在这方面获得的结果至今仍被认为是该理论中最精妙的结果之一）。“群的线性表示”的最初想法也追溯到伽罗瓦†，这一事实清楚地表明他拥有两个群结构的同构概念，而不依赖于它们的“实现”。

然而，尽管高斯和伽罗瓦的巧妙方法无可争辩地引导他们形成了对合成律概念的非常宽广的理解，但他们并未特别在这方面发展他们的思想，他们的工作对抽象代数的演进也没有产生直接的影响。‡ 朝向抽象的最清晰进展是在第三个方向上取得的：在反思虚数的性质之后（虚数的几何表示在 19 世纪初引发了大量的工作），英国学派的代数学家们率先在 1830 年至 1850 年间分离出合成律的抽象概念，并立即通过将这一概念应用于大量新的数学实体而拓宽了代数的领域：布尔那里的逻辑代数（参见《集合论》第 IV 章的历史注记），哈密顿那里的向量和四元数，凯莱那里的一般超复数系统、矩阵与非结合律（{[}10{]}，第 I 卷，第 127 页和第 301 页，第 II 卷，第 185 页和第 475 页）。在欧洲大陆也独立地进行了平行的演进，特别是在向量演算（莫比乌斯、贝拉维蒂斯）、线性代数和超复数系统（格拉斯曼）方面，我们将在《代数》第 III 章的历史注记§中更详细地讨论这些内容。

从所有这些在 19 世纪上半叶使代数焕发活力的原创而富有成果的思想的激荡中，代数彻底地焕然一新。此前，方法和结果都围绕着代数方程（或数论中的丢番图方程）求解这一中心问题而展开：“代数”，塞雷特在其《高等代数教程》的引言中说道 {[}12{]} ，“严格说来，是方程的分析”。1850 年之后，尽管代数方面的论著在一段时间内仍然给予方程理论以优先地位，但新的研究工作不再受制于对数值方程求解的直接应用的关注，而是越来越多地朝向今天被认为是代数的本质问题的方向，即为研究而研究代数结构。

这些工作相当整齐地分为三大主流，它们分别延续了我们上面分析过的三个思想运动，并且在19世纪最后几年之前一直平行发展，彼此之间没有显著的影响。†

这些工作主流中的第一个源于 19 世纪德国学派（狄利克雷、库默尔、克罗内克、戴德金、希尔伯特）在建立代数数论方面的工作，其发端于高斯的研究，高斯开创了这类数的第一种研究，即数 \(a + bi\) （ \(a\) 和 \(b\) 为有理数）的研究。我们在此不追踪这一理论的发展（参见《交换代数》第七章的历史注记）：就我们的目的而言，只需提及由此产生的抽象代数概念。从高斯的最初后继者起，域（代数数域）的概念就是这一主题所有著作的基础（正如它对阿贝尔和伽罗瓦关于代数方程的研究一样）；当戴德金和韦伯 {[}13{]} 以代数数理论为模型建立了单变量代数函数理论时，它的应用范围得到了扩大。戴德金 {[}14{]} 还提出了理想的概念，这一概念为元素集合之间的合成律提供了新的范例；他和克罗内克共同推动了交换群和模在代数域理论中日益重要的作用；我们将在《代数》第三、五、七章的历史注记中再次讨论这一点。

关于线性代数与超复数系统的发展，我们还将参考《代数》第 III 章和第 VIII 章的历史注记；这一发展在 19 世纪末至 20 世纪初并未引入任何新的代数概念，它在英国（西尔维斯特、W. 克利福德）和美国（B. 皮尔斯和 C. S. 皮尔斯、迪克森、韦德伯恩）沿着哈密顿和凯莱所开辟的道路进行，同时在德国（魏尔斯特拉斯、戴德金、弗罗贝尼乌斯、莫利恩）和法国（拉盖尔、E. 嘉当）独立于盎格鲁-撒克逊学派，并采用了截然不同的方法。

就群论而言，它最初主要是作为有限置换群的理论发展起来的，这发生在伽罗瓦的著作出版之后，并通过塞雷特的著作 {[}12{]} 以及尤其是 C. 若尔当的伟大著作《置换论》 {[}15{]} 的传播而实现。若尔当在那里极大地改进了其前辈关于置换群特殊性质（传递性、本原性等）的研究成果，获得了大部分至今未被超越的结果；他还对一类非常重要的特殊群——线性群及其子群——进行了同样深入的研究；此外，正是他引入了从一个群到另一个群的同态这一基本概念，并且（稍后）引入了商群的概念，还证明了被称为“若尔当–赫尔德定理”的定理的一部分。† 最后，若尔当是第一个研究无限群的人 {[}16{]} ，这些无限群后来由 S. 李一方面以及 F. 克莱因和 H. 庞加莱另一方面在几年后沿着两个不同方向进行了相当大的发展。

与此同时，群中本质的东西，即其合成律而非构成群的对象的性质，逐渐被认识到（例如参见 {[}10{]}，第 II 卷，第 123 页和第 131 页以及 {[}14{]}，第 III 卷，第 439 页）。然而，即使是对有限抽象群的研究工作，在很长一段时间内也被视为置换群的研究，直到大约 1880 年，有限群的独立理论才开始有意识地发展起来。我们无法进一步追溯这一理论的历史，本专著中对此仅作了非常肤浅的涉及；我们只提及两个至今在研究有限群中仍最常用的工具，它们都追溯到 19 世纪：关于 p 群的西罗定理‡（其年代为 1872 年 {[}17{]} ）以及该世纪末由弗罗贝尼乌斯创立的特征标理论 {[}19{]} 。我们建议希望深入研究有限群理论及其引发的许多难题的读者，参阅伯恩赛德 {[}20{]} 、施派泽 {[}23{]} 、扎森豪斯 {[}24{]} 和戈伦斯坦 {[}27{]} 的专著。

大约在 1880 年，群表现的系统研究也开始了；此前仅遇到过特定群的表现，例如交错群 \(\mathfrak{A}_5\) 在哈密顿的一篇著作中 {[}9 bis{]} 的表现，或线性微分方程和黎曼曲面的单值群（施瓦茨、克莱因、施莱夫利）的表现。W. 迪克是第一个 {[}18{]} 定义（但未给它命名）由有限个生成元生成的自由群的人，但他对自由群本身的兴趣不如把它作为“万有”对象那样大，后者使得可以精确定义“由生成元和关系”给出的群。迪克发展了凯莱首先提出、同时又明显受到上述黎曼学派著作影响的思想，描述了具有给定表现的群的一种解释，其中每个生成元由关于相切或相交圆的两个反演之积来表示（另见伯恩赛德 {[}20{]} ，第十八章）。稍后，在庞加莱及其后继者发展了自守函数理论之后，庞加莱又引入了代数拓扑的工具，对基本群的初步研究便与群表现的研究（德恩、蒂策）并驾齐驱，两种理论相互支持。此外，正是拓扑学家 J. 尼尔森在 1924 年对自由群性质的第一个深入研究中引入了“自由群”这一术语 {[}21{]} ；紧接着，E. 阿廷（始终参照拓扑学中的问题）引入了群的自由积的概念，而 O. 施赖尔则更一般地定义了带融合子群的自由积 {[}22{]} （参见 I，§ 7，第 3 号，习题 20）。在此我们同样无法详述这一方向后续发展的历史，请读者参阅著作 {[}26{]} 以获取更多细节。

自 19 世纪末以来，群的概念（以及与之密切相关的“不变量”概念）在分析学、几何学、力学和理论物理学中所取得的非凡成功，已无需再多言。这一概念以及与之相关的代数概念（带算子的群、环、理想、模）同样侵入到此前看似毫无关联的代数学的各个部分，这构成了我们在此追溯的、并以我们上文所追踪的三个分支的综合而告终的最后一个发展时期的显著标志。这一统一主要是 1900 年至 1930 年间德国学派的工作成果：由戴德金和希尔伯特在 19 世纪最后几年开创的代数学公理化工作，由 E. 施泰尼茨有力地继续推进，随后自 1920 年起，在 E. 阿廷、E. 诺特及其学派代数学家（哈塞、克鲁尔、O. 施赖尔、范德瓦尔登）的推动下得到大力发展。范德瓦尔登的专著 {[}25{]} 于 1930 年出版，首次将这些工作汇集于一部统一的论述之中，为二十多年间代数学领域的许多研究工作开辟了道路并起到了指导作用。

